\section{Conclusion}

Geographic gaming is bounded and detectable:
\begin{enumerate}
  \item \textbf{Resolution gaming} produces at most $0.18$ seat shift nationally
    from varying the F.3 threshold across its plausible range.
  \item \textbf{Adjacency manipulation} (removing 50 cross-partisan edges in NC)
    produces a $0.15$ seat shift.
  \item \textbf{Combined maximum}: even if an adversary simultaneously chose the
    most favourable resolution threshold and manipulated the adjacency graph,
    the national partisan shift would be $< 0.35$ seats --- within sampling noise.
  \item \textbf{Detectability}: both vectors are detectable via the audit chain.
    Resolution is recorded in the manifest; adjacency is hashed against the
    standard TIGER/Line construction.
\end{enumerate}

The key design principle: geographic definition choices are a pre-registration
problem, not a manipulation-resistance problem.
If the resolution rule and adjacency definition are specified and pre-registered
before the Census data is released (as the DIA requires), there is no opportunity
for post-data geographic gaming.
The audit chain detects any deviation from the pre-registered specification.
